\documentclass[14pt,a4paper]{extreport}
% preamble.tex — спільна преамбула для всіх лабораторних робіт. Підключається одразу після \documentclass у кожному lab*/main.tex.
%
% ЗБІРКА: LuaLaTeX (не pdfLaTeX). Причина: пакет listings під pdfLaTeX не читає UTF-8 і зупиняє збірку на кириличних коментарях у лістингах програмного коду. LuaLaTeX обробляє UTF-8 нативно.
%
% ШРИФТ: Liberation Serif — метрично сумісний з Times New Roman і містить повний набір кириличних гліфів (TeX Gyre Termes кирилиці не має, через що український текст не відображався взагалі).

\usepackage{fontspec}
\setmainfont{Liberation Serif}
\setmonofont{DejaVu Sans Mono}[Scale=0.82]

\usepackage{polyglossia}
\setmainlanguage{ukrainian}
\setotherlanguage{english}

\usepackage{microtype}

% Поля, інтервал, відступи й заголовки нижче взяті з шаблону НУЛП, за яким написана бакалаврська робота (../bachelor-thesis/thesis/main.tex). Той самий формат приймає В.М. Шиманський і для магістерських, тож усі звіти тримаємо в ньому, а не в тому, що переказували на консультації.
\usepackage[left=25mm,right=15mm,top=15mm,bottom=23mm,footskip=1.5mm]{geometry}
\usepackage{setspace}
% Саме \setstretch{1.5}, а не \onehalfspacing: останній дає близько 1.25.
\setstretch{1.5}

\usepackage{graphicx}
\usepackage{float}
\usepackage{booktabs}
\usepackage{array}
\usepackage{tabularx}
\usepackage{longtable}
\usepackage{enumitem}
% labelindent вирівнює маркер списку з абзацним відступом — так у шаблоні НУЛП.
\setlist{nosep}
\setlist[itemize]{label=\textbullet,leftmargin=*,labelindent=1.27cm}
\setlist[enumerate]{leftmargin=*,labelindent=1.27cm}
\usepackage{indentfirst}
\usepackage{amsmath}

\usepackage{hyperref}
\hypersetup{colorlinks=true,linkcolor=black,urlcolor=blue,citecolor=black}

% xurl дозволяє \url розривати адресу в будь-якому місці, а не лише після скісної риски. Без нього довгий сегмент шляху (напр. «scopus---podrobnaya- instruktsiya-dlya-uchenogo») виходить за межі сторінки, бо дефіс для hyperref не є точкою розриву. Вантажиться після hyperref — перевизначає його \UrlBreaks.
\usepackage{xurl}

% ---------------------------- Список використаної літератури ----------------------------
% biblatex + biber зі стилем IEEE: посилання «[1]» у тексті, записи — за порядком першої згадки. Обрано IEEE, а не gost-numeric (ДСТУ 8302:2015), — так вирішив автор звітів після порівняння обох варіантів.
% Кожен звіт із джерелами сам підключає свій references.bib через \addbibresource і друкує список \printbibliography[heading=none] під власним \section. Звіт без джерел від завантаженого biblatex нічого не втрачає.
% Оточення списку — стандартне для стилю: мітка від лівого поля, текст опису вирівняний після неї. Не підганяємо його під абзацний відступ нумерованих переліків тексту (labelindent у \setlist вище) — список літератури є окремим розділом зі своїм стандартним оформленням.
\usepackage[backend=biber,style=ieee]{biblatex}
% biblatex за замовчуванням додає розтяжний проміжок після «/ . : ? =» в адресах, і з вимкненими переносами рядок вирівнюється розрідженням самої адреси («https : / / www . scopus»). Рвати адресу й так дозволяє xurl, тож проміжок лише шкодить.
\setlength{\biburlbigskip}{0mu}
% Без додаткового інтервалу між описами: розділ має йти суцільним текстом, як решта звіту. Так само в шаблоні бакалаврської.
\setlength{\bibitemsep}{0pt}
\setlength{\bibparsep}{0pt}

% ---------------------------- Лістинги програмного коду ---------------------------- Використано minted (через Pygments), а не listings: listings не вміє коректно обробляти багатобайтові символи й «висмикує» кириличні коментарі з коду, зліплюючи їх в один рядок. Потребує прапорця -shell-escape.
\usepackage{xcolor}
\usepackage[newfloat]{minted}
% labelsep=period дає «Рис. 2.1. Назва» замість «Рис. 2: Назва» — формат НУЛП.
\usepackage[labelsep=period]{caption}

% baselinestretch=1 обов'язковий: без нього лістинг успадковує півторашку основного тексту, і довгий фрагмент коду перестає вміщатися у сторінку. Усередині [H]-флоата це не помилка збірки, а нескінченний цикл порожніх сторінок — lab1 так намолотив 161 тисячу сторінок.
\setminted{
  baselinestretch=1,
  fontsize=\footnotesize,
  linenos=true,
  frame=single,
  framesep=4pt,
  breaklines=true,
  breakanywhere=true,
  tabsize=4,
  numbersep=7pt,
  bgcolor=gray!8,
}
\setminted[python]{python3=true}

\SetupFloatingEnvironment{listing}{name=Лістинг,listname=Перелік лістингів}

% ---------------------------- Нумерація та оформлення заголовків ----------------------------
\usepackage{titlesec}

\setcounter{secnumdepth}{3}
\renewcommand{\thesection}{\arabic{section}}
\renewcommand{\thesubsection}{\arabic{section}.\arabic{subsection}}
\renewcommand{\thesubsubsection}{\arabic{section}.\arabic{subsection}.\arabic{subsubsection}}

% Розділ по центру, 16 пт, великими літерами; слово «РОЗДІЛ» опущене, бо звіт не дисертація. Номер лишається: \numberwithin прив'язує до нього нумерацію рисунків («Рис. 2.1»), і без нумерованих розділів вона б зламалася. \MakeUppercase четвертим аргументом отримує текст заголовка, тому в джерелі заголовки пишуться звичайним регістром.
\titleformat{\section}[block]
  {\centering\bfseries\fontsize{16}{19.2}\selectfont}
  {\thesection.}{0.3em}{\MakeUppercase}
% Підрозділи — зліва з абзацного відступу, звичайним регістром.
\titleformat{\subsection}[block]{\bfseries}{\thesubsection.}{1em}{}
\titleformat{\subsubsection}[block]{\bfseries\itshape}{\thesubsubsection.}{1em}{}

% Незіркові \titlespacing — обов'язкові. Без них titlesec застосовує зірковий варіант, який прибирає абзацний відступ у першого абзацу після заголовка й перекриває indentfirst. Значення — типові для report.
\titlespacing{\section}{0pt}{20pt}{6pt}
\titlespacing{\subsection}{1.27cm}{6pt}{0pt}
\titlespacing{\subsubsection}{1.27cm}{6pt}{0pt}

\setlength{\parindent}{1.27cm}
\setlength{\parskip}{0pt}
\renewcommand{\arraystretch}{1.3}

% Переноси слів вимкнені — так у шаблоні НУЛП. Натомість абзац розганяється міжслівними пробілами, тому tolerance і emergencystretch підняті: без них рядки виходили б за праве поле.
\hyphenpenalty=10000
\exhyphenpenalty=10000
\tolerance=9999
\setlength{\emergencystretch}{3em}

% ---------------------------- Підписи до рисунків і таблиць ----------------------------
\captionsetup[figure]{name=Рис.,justification=centering,font={bf,it}}
\captionsetup[table]{name=Таблиця,justification=raggedleft,singlelinecheck=false,labelfont=bf,textfont=normalfont}

% Рисунки, таблиці й формули нумеруються в межах розділу: «Рис. 2.1», а не «Рис. 5». Так у шаблоні НУЛП, і так само виглядали приклади на консультації.
\numberwithin{equation}{section}
\numberwithin{figure}{section}
\numberwithin{table}{section}

% ---------------------------- Колонтитули ----------------------------
\usepackage{fancyhdr}
\pagestyle{fancy}
\fancyhf{}
\fancyfoot[R]{\thepage}
\renewcommand{\headrulewidth}{0pt}
\fancypagestyle{plain}{%
  \fancyhf{}
  \fancyfoot[R]{\thepage}
  \renewcommand{\headrulewidth}{0pt}
}

% titlepage.tex — титульний аркуш, спільний для всіх дисциплін.
% Потребує визначених у subject.tex команд:
%   \SubjectName, \LecturerName, \StudentName, \StudentGroup, \AcademicYear
%
% Використання у тексті: \MakeLabTitle{1}{Тема роботи}
%
% Рід роботи. Дисципліни з практичними роботами перевизначають його після
% % preamble.tex — спільна преамбула для всіх лабораторних робіт. Підключається одразу після \documentclass у кожному lab*/main.tex.
%
% ЗБІРКА: LuaLaTeX (не pdfLaTeX). Причина: пакет listings під pdfLaTeX не читає UTF-8 і зупиняє збірку на кириличних коментарях у лістингах програмного коду. LuaLaTeX обробляє UTF-8 нативно.
%
% ШРИФТ: Liberation Serif — метрично сумісний з Times New Roman і містить повний набір кириличних гліфів (TeX Gyre Termes кирилиці не має, через що український текст не відображався взагалі).

\usepackage{fontspec}
\setmainfont{Liberation Serif}
\setmonofont{DejaVu Sans Mono}[Scale=0.82]

\usepackage{polyglossia}
\setmainlanguage{ukrainian}
\setotherlanguage{english}

\usepackage{microtype}

% Поля, інтервал, відступи й заголовки нижче взяті з шаблону НУЛП, за яким написана бакалаврська робота (../bachelor-thesis/thesis/main.tex). Той самий формат приймає В.М. Шиманський і для магістерських, тож усі звіти тримаємо в ньому, а не в тому, що переказували на консультації.
\usepackage[left=25mm,right=15mm,top=15mm,bottom=23mm,footskip=1.5mm]{geometry}
\usepackage{setspace}
% Саме \setstretch{1.5}, а не \onehalfspacing: останній дає близько 1.25.
\setstretch{1.5}

\usepackage{graphicx}
\usepackage{float}
\usepackage{booktabs}
\usepackage{array}
\usepackage{tabularx}
\usepackage{longtable}
\usepackage{enumitem}
% labelindent вирівнює маркер списку з абзацним відступом — так у шаблоні НУЛП.
\setlist{nosep}
\setlist[itemize]{label=\textbullet,leftmargin=*,labelindent=1.27cm}
\setlist[enumerate]{leftmargin=*,labelindent=1.27cm}
\usepackage{indentfirst}
\usepackage{amsmath}

\usepackage{hyperref}
\hypersetup{colorlinks=true,linkcolor=black,urlcolor=blue,citecolor=black}

% xurl дозволяє \url розривати адресу в будь-якому місці, а не лише після скісної риски. Без нього довгий сегмент шляху (напр. «scopus---podrobnaya- instruktsiya-dlya-uchenogo») виходить за межі сторінки, бо дефіс для hyperref не є точкою розриву. Вантажиться після hyperref — перевизначає його \UrlBreaks.
\usepackage{xurl}

% ---------------------------- Список використаної літератури ----------------------------
% biblatex + biber зі стилем IEEE: посилання «[1]» у тексті, записи — за порядком першої згадки. Обрано IEEE, а не gost-numeric (ДСТУ 8302:2015), — так вирішив автор звітів після порівняння обох варіантів.
% Кожен звіт із джерелами сам підключає свій references.bib через \addbibresource і друкує список \printbibliography[heading=none] під власним \section. Звіт без джерел від завантаженого biblatex нічого не втрачає.
% Оточення списку — стандартне для стилю: мітка від лівого поля, текст опису вирівняний після неї. Не підганяємо його під абзацний відступ нумерованих переліків тексту (labelindent у \setlist вище) — список літератури є окремим розділом зі своїм стандартним оформленням.
\usepackage[backend=biber,style=ieee]{biblatex}
% biblatex за замовчуванням додає розтяжний проміжок після «/ . : ? =» в адресах, і з вимкненими переносами рядок вирівнюється розрідженням самої адреси («https : / / www . scopus»). Рвати адресу й так дозволяє xurl, тож проміжок лише шкодить.
\setlength{\biburlbigskip}{0mu}
% Без додаткового інтервалу між описами: розділ має йти суцільним текстом, як решта звіту. Так само в шаблоні бакалаврської.
\setlength{\bibitemsep}{0pt}
\setlength{\bibparsep}{0pt}

% ---------------------------- Лістинги програмного коду ---------------------------- Використано minted (через Pygments), а не listings: listings не вміє коректно обробляти багатобайтові символи й «висмикує» кириличні коментарі з коду, зліплюючи їх в один рядок. Потребує прапорця -shell-escape.
\usepackage{xcolor}
\usepackage[newfloat]{minted}
% labelsep=period дає «Рис. 2.1. Назва» замість «Рис. 2: Назва» — формат НУЛП.
\usepackage[labelsep=period]{caption}

% baselinestretch=1 обов'язковий: без нього лістинг успадковує півторашку основного тексту, і довгий фрагмент коду перестає вміщатися у сторінку. Усередині [H]-флоата це не помилка збірки, а нескінченний цикл порожніх сторінок — lab1 так намолотив 161 тисячу сторінок.
\setminted{
  baselinestretch=1,
  fontsize=\footnotesize,
  linenos=true,
  frame=single,
  framesep=4pt,
  breaklines=true,
  breakanywhere=true,
  tabsize=4,
  numbersep=7pt,
  bgcolor=gray!8,
}
\setminted[python]{python3=true}

\SetupFloatingEnvironment{listing}{name=Лістинг,listname=Перелік лістингів}

% ---------------------------- Нумерація та оформлення заголовків ----------------------------
\usepackage{titlesec}

\setcounter{secnumdepth}{3}
\renewcommand{\thesection}{\arabic{section}}
\renewcommand{\thesubsection}{\arabic{section}.\arabic{subsection}}
\renewcommand{\thesubsubsection}{\arabic{section}.\arabic{subsection}.\arabic{subsubsection}}

% Розділ по центру, 16 пт, великими літерами; слово «РОЗДІЛ» опущене, бо звіт не дисертація. Номер лишається: \numberwithin прив'язує до нього нумерацію рисунків («Рис. 2.1»), і без нумерованих розділів вона б зламалася. \MakeUppercase четвертим аргументом отримує текст заголовка, тому в джерелі заголовки пишуться звичайним регістром.
\titleformat{\section}[block]
  {\centering\bfseries\fontsize{16}{19.2}\selectfont}
  {\thesection.}{0.3em}{\MakeUppercase}
% Підрозділи — зліва з абзацного відступу, звичайним регістром.
\titleformat{\subsection}[block]{\bfseries}{\thesubsection.}{1em}{}
\titleformat{\subsubsection}[block]{\bfseries\itshape}{\thesubsubsection.}{1em}{}

% Незіркові \titlespacing — обов'язкові. Без них titlesec застосовує зірковий варіант, який прибирає абзацний відступ у першого абзацу після заголовка й перекриває indentfirst. Значення — типові для report.
\titlespacing{\section}{0pt}{20pt}{6pt}
\titlespacing{\subsection}{1.27cm}{6pt}{0pt}
\titlespacing{\subsubsection}{1.27cm}{6pt}{0pt}

\setlength{\parindent}{1.27cm}
\setlength{\parskip}{0pt}
\renewcommand{\arraystretch}{1.3}

% Переноси слів вимкнені — так у шаблоні НУЛП. Натомість абзац розганяється міжслівними пробілами, тому tolerance і emergencystretch підняті: без них рядки виходили б за праве поле.
\hyphenpenalty=10000
\exhyphenpenalty=10000
\tolerance=9999
\setlength{\emergencystretch}{3em}

% ---------------------------- Підписи до рисунків і таблиць ----------------------------
\captionsetup[figure]{name=Рис.,justification=centering,font={bf,it}}
\captionsetup[table]{name=Таблиця,justification=raggedleft,singlelinecheck=false,labelfont=bf,textfont=normalfont}

% Рисунки, таблиці й формули нумеруються в межах розділу: «Рис. 2.1», а не «Рис. 5». Так у шаблоні НУЛП, і так само виглядали приклади на консультації.
\numberwithin{equation}{section}
\numberwithin{figure}{section}
\numberwithin{table}{section}

% ---------------------------- Колонтитули ----------------------------
\usepackage{fancyhdr}
\pagestyle{fancy}
\fancyhf{}
\fancyfoot[R]{\thepage}
\renewcommand{\headrulewidth}{0pt}
\fancypagestyle{plain}{%
  \fancyhf{}
  \fancyfoot[R]{\thepage}
  \renewcommand{\headrulewidth}{0pt}
}

\input{../../../../titlepage}
: \renewcommand{\WorkKind}{ПРАКТИЧНА РОБОТА}.
\newcommand{\WorkKind}{ЛАБОРАТОРНА РОБОТА}

\newcommand{\MakeLabTitle}[2]{%
\begin{titlepage}
% Титулка — не суцільний текст, а фіксована верстка з явних \\[...] проміжків,
% тож міжрядковий інтервал основного тексту (\setstretch{1.5} у preamble.tex)
% до неї не застосовний: він розтягував кожен із ~15 рядків на 20 %, і блок
% «Виконав» разом із «Львів --- рік» переповнював сторінку. Інтервал локальний
% для оточення titlepage, тому на решту документа це не впливає.
\setstretch{1}
\begin{center}
  {\large\bfseries МІНІСТЕРСТВО ОСВІТИ І НАУКИ УКРАЇНИ}\\[6pt]
  {\large\bfseries Національний університет «Львівська політехніка»}\\[6pt]
  {\large\bfseries Інститут комп'ютерних наук та інформаційних технологій}\\[6pt]
  {\large\bfseries Кафедра «Системи штучного інтелекту»}\\[1.2cm]

  \includegraphics[width=5cm]{../../../../logo}\\[1.2cm]

  {\Large\bfseries \WorkKind{} №#1}\\[0.5cm]
  з дисципліни\\[4pt]
  «\SubjectName»\\[0.6cm]
  на тему\\[4pt]
  \textit{#2}
\end{center}

% Відступ до блоку «Виконав» гнучкий, а не фіксований (було 2 cm). Фіксований
% разом із довгим родом роботи чи темою на два рядки переповнював сторінку, і
% «Львів --- рік» зі свого \vfill з'їжджав на окрему другу сторінку. \vfill
% стискається до нуля, коли місця бракує, тож переповнення стало неможливим.
\vfill

\begin{flushright}
  {\bfseries Виконав:}\\
  студент групи \StudentGroup\\
  \StudentName\\[1cm]
  {\bfseries Викладач:}\\
  \LecturerName\\
\end{flushright}

\vfill
\begin{center}
  Львів --- \AcademicYear
\end{center}
\end{titlepage}
\clearpage
}


% Метадані дисципліни. Використовуються титульним аркушем (titlepage.tex).

\newcommand{\SubjectName}{Штучний інтелект та інноваційні інформаційні технології}
% Два рядки: посада з науковим ступенем, далі прізвище. `\\` розриває рядок
% усередині flushright-блоку титульного аркуша.
\newcommand{\LecturerName}{доц. каф. СШІ к.т.н.,\\ Думин І.Б.}
\newcommand{\StudentName}{Галич Максим}
\newcommand{\StudentGroup}{КНСШ-13}
\newcommand{\AcademicYear}{2026}

% Рисунки — згенеровані артефакти, тому лежать не поряд із текстом, а в .cache/ (див. ARTIFACTS у src/aiit/config.py).
\graphicspath{{../../../../.cache/reports-artifacts/1th term 2026 Autumn/Штучний інтелект та інноваційні інформаційні технології/}}

\begin{document}

\MakeLabTitle{3}{Використання еволюційних алгоритмів}

\textbf{Мета роботи.} Ознайомитися з принципами еволюційних алгоритмів як одного з напрямів штучного інтелекту, отримати практичні навички їх використання для розв'язання оптимізаційних задач.

\section{Постановка задачі}

Виконано \textbf{варіант 1}: власна програмна реалізація генетичного алгоритму (ГА), який підбирає гіперпараметри моделі машинного навчання.

Модель і предметну область узято з лабораторних робіт №1--2: повнозв'язна нейронна мережа у PyTorch розв'язує задачу регресії на наборі даних \textbf{California Housing} (20\,640 записів, 8 ознак, цільова змінна --- медіанна вартість будинку в сотнях тисяч доларів). Поділ вибірки той самий: 60/20/20 зі стратифікацією, тобто 12\,384 записи для навчання, 4\,128 для валідації і 4\,128 для тестування.

Такий вибір має практичну мотивацію. У лабораторній роботі №2 гіперпараметри задавалися вручну, і регуляризація, додана «про всяк випадок», погіршила результат на 22~\%. Еволюційний пошук дає змогу замінити ручний перебір систематичним і перевірити, наскільки далеким від оптимуму був ручний вибір.

\textbf{Функція пристосованості} --- середньоквадратична помилка (RMSE) моделі на валідаційній вибірці після 20 епох навчання. Її потрібно мінімізувати. Тестова вибірка в пошуку не бере участі: на ній вимірюються лише фінальні конфігурації після завершення пошуку, інакше тестова RMSE стала б ще однією величиною, під яку підганяється модель.

\section{Простір пошуку та кодування розв'язків}

Кожен розв'язок (особина) --- це набір із восьми гіперпараметрів, тобто вісім генів. Кожен ген набуває одного зі скінченного переліку значень (табл.~\ref{tab:space}), тому хромосома кодується як кортеж із восьми цілих чисел --- індексів значень у цих переліках. Таке кодування дає змогу застосовувати однакові оператори до генів різної природи: числових (швидкість навчання) і категоріальних (функція активації).

\begin{table}[H]
\caption{Гени хромосоми та їх допустимі значення}
\label{tab:space}
\begin{tabularx}{\linewidth}{|l|X|l|}
\hline
\textbf{Ген} & \textbf{Значення} & \textbf{Тип} \\
\hline
Кількість прихованих шарів & 1, 2, 3 & впорядкований \\
\hline
Ширина першого шару & 32, 64, 128, 256 & впорядкований \\
\hline
Функція активації & ReLU, GELU, LeakyReLU & категоріальний \\
\hline
Dropout & 0; 0,1; 0,2; 0,3 & впорядкований \\
\hline
Batch Normalization & ні, так & категоріальний \\
\hline
Weight decay (L2) & 0; $10^{-5}$; $10^{-4}$; $10^{-3}$ & впорядкований \\
\hline
Швидкість навчання Adam & $3 \cdot 10^{-4}$; $10^{-3}$; $3 \cdot 10^{-3}$; $10^{-2}$ & впорядкований \\
\hline
Розмір батча & 64, 128, 256, 512 & впорядкований \\
\hline
\end{tabularx}
\end{table}

Кожен наступний прихований шар удвічі вужчий за попередній, тож ширина першого шару і кількість шарів однозначно задають архітектуру: наприклад, 2 шари з шириною 64 дають мережу $64 \to 32$. Простір пошуку містить $3 \cdot 4 \cdot 3 \cdot 4 \cdot 2 \cdot 4 \cdot 4 \cdot 4 = 18\,432$ конфігурації. Повний перебір потребував би 18\,432 навчань, тоді як генетичний алгоритм обмежено бюджетом приблизно в 120 навчань, тобто менш ніж 1~\% простору.

Найкраща модель лабораторної роботи №2 (baseline-relu) є однією з точок цього простору: 2 шари, ширина 64, ReLU, без Dropout, BatchNorm і weight decay, швидкість навчання $10^{-3}$, батч 256. Вона слугує точкою відліку для оцінки результатів пошуку.

\begin{listing}[H]
\inputminted[firstline=56,lastline=66]{python}{../src/aiit/run_lab3.py}
\caption{Простір пошуку (\texttt{src/aiit/run\_lab3.py})}
\label{lst:space}
\end{listing}

\section{Генетичний алгоритм}

\subsection{Схема алгоритму}

Алгоритм реалізовано самостійно, без сторонніх бібліотек еволюційних обчислень, у модулі \texttt{src/aiit/evolution.py}. Один запуск складається з таких кроків.

\begin{enumerate}
\item \textbf{Ініціалізація.} Генерується початкова популяція з 12 різних випадкових хромосом.
\item \textbf{Оцінювання.} Для кожної нової хромосоми навчається відповідна мережа і обчислюється її валідаційна RMSE.
\item \textbf{Елітизм.} Дві найкращі особини переходять у наступне покоління без змін. Тому найкраща пристосованість популяції ніколи не погіршується.
\item \textbf{Відбір.} Батьків для решти місць обирає турнірний відбір: з популяції випадково беруться три особини, і перемагає та, що має найменшу RMSE.
\item \textbf{Схрещування.} З імовірністю 0,9 нащадок утворюється рівномірним схрещуванням двох батьків, інакше він є копією першого батька.
\item \textbf{Мутація.} Кожен ген нащадка мутує з імовірністю $1/8$.
\item \textbf{Відкидання дублікатів.} Якщо нащадок збігається з особиною, яка вже є в новому поколінні, він мутує далі, доки не стане унікальним.
\item \textbf{Контроль виродження.} Якщо різноманітність нащадків опускається нижче порогу 0,3, три останні нащадки замінюються випадковими хромосомами --- «іммігрантами».
\item Кроки 2--8 повторюються, доки не буде сформовано 15 поколінь, включно з початковим.
\end{enumerate}

\subsection{Еволюційні оператори}

\textbf{Турнірний відбір} обрано замість рулеткового, бо він залежить лише від порядку особин за пристосованістю, а не від абсолютних значень RMSE. Одна мережа, що розійшлася під час навчання і отримала RMSE $= +\infty$, не спотворює ймовірностей відбору решти особин, а одна дуже вдала особина не захоплює всю популяцію за одне покоління. Розмір турніру 3 дає помірний тиск відбору: для популяції з 12 особин найгірші особини майже ніколи не стають батьками, але середні мають реальний шанс.

\textbf{Рівномірне схрещування} обрано замість одноточкового, бо порядок гіперпараметрів у хромосомі довільний. Одноточкове схрещування успадковувало б сусідні гени разом, тобто штучно пов'язувало б, наприклад, Dropout із Batch Normalization лише тому, що вони стоять поруч у кортежі.

\textbf{Мутація} враховує тип гена. Для впорядкованих генів вона змінює значення на сусіднє (наприклад, швидкість навчання $10^{-3}$ стає $3 \cdot 10^{-4}$ або $3 \cdot 10^{-3}$): так вдале значення уточнюється, а не втрачається. Для категоріальних генів, де поняття «сусіднього» значення немає, обирається будь-яке інше значення рівноймовірно. Імовірність мутації $1/n$, де $n = 8$ --- кількість генів, є стандартним вибором: у середньому мутує один ген на нащадка.

% Без оточення listing: лістинг вищий за сторінку, а плаваючий об'єкт розривати не можна — [H] не має куди його покласти, і збірка зациклюється. Звичайний \inputminted переноситься між сторінками, \captionof дає ту саму нумерацію.
\inputminted[firstline=192,lastline=241]{python}{../src/aiit/evolution.py}
\captionof{listing}{Турнірний відбір, рівномірне схрещування та мутація (\texttt{src/aiit/evolution.py})}
\label{lst:operators}

\subsection{Захист від виродження популяції}

Виродження --- це стан, коли всі особини стають майже однаковими. Тоді схрещування не створює нічого нового, бо обоє батьків несуть ті самі гени, і пошук зупиняється. Для малої популяції (12 особин) з елітизмом цей ризик особливо високий: елітні особини постійно перемагають у турнірах, і їхні копії заповнюють популяцію.

Реалізовано два механізми захисту. Перший --- \textbf{відкидання дублікатів}: нащадок, що збігається з уже сформованою особиною нового покоління, мутує далі. Другий --- \textbf{іммігранти}: якщо різноманітність падає нижче порогу, частину нащадків замінюють випадковими хромосомами. Еліту іммігранти не заміщують.

Різноманітність популяції вимірюється як середня частка генів, якими відрізняються дві особини, по всіх парах особин. Значення 1 означає, що будь-які дві особини відрізняються всіма генами, значення 0 --- що популяція складається з копій однієї хромосоми. Поріг 0,3 означає, що дві особини в середньому відрізняються менш ніж 2,4 генами з 8.

\inputminted[firstline=271,lastline=310]{python}{../src/aiit/evolution.py}
\captionof{listing}{Формування наступного покоління (\texttt{src/aiit/evolution.py})}
\label{lst:nextgen}

\subsection{Параметри алгоритму}

\begin{table}[H]
\caption{Параметри генетичного алгоритму та їх обґрунтування}
\label{tab:gaparams}
\begin{tabularx}{\linewidth}{|l|c|X|}
\hline
\textbf{Параметр} & \textbf{Значення} & \textbf{Обґрунтування} \\
\hline
Розмір популяції & 12 & Два пакети по 6 паралельних навчань на покоління; бюджет $\approx 120$ навчань \\
\hline
Кількість поколінь & 15 & Пробний запуск із 10 поколіннями покращував результат до 8-го покоління; 15 поколінь дають запас, щоб побачити вихід на плато \\
\hline
Еліта & 2 & Гарантує монотонність найкращого результату; більша еліта зменшує частку нових особин \\
\hline
Розмір турніру & 3 & Помірний тиск відбору (25~\% популяції) \\
\hline
Імовірність схрещування & 0,9 & Більшість нащадків поєднує гени двох батьків \\
\hline
Імовірність мутації гена & $1/8$ & У середньому один змінений ген на нащадка \\
\hline
Поріг різноманітності & 0,3 & Менше 2,4 відмінних генів із 8 у середньому \\
\hline
Кількість іммігрантів & 3 & 25~\% популяції; еліта не заміщується \\
\hline
\end{tabularx}
\end{table}

\inputminted[firstline=89,lastline=103]{python}{../src/aiit/run_lab3.py}
\captionof{listing}{Параметри ГА та «наївного» ГА без захисту від виродження (\texttt{src/aiit/run\_lab3.py})}
\label{lst:gaconfig}

\section{Обчислення пристосованості}

Щоб обчислити пристосованість хромосоми, її декодують у набір гіперпараметрів, будують за ним мережу, навчають її 20 епох оптимізатором Adam і обчислюють RMSE на валідаційній вибірці.

Навчання повністю детерміноване: ініціалізацію ваг і порядок батчів задано фіксованим зерном 42. Тому одна й та сама хромосома завжди отримує ту саму пристосованість, і кожну унікальну хромосому достатньо навчити лише один раз. Повторні оцінки, зокрема елітних особин, які переходять у наступне покоління, беруться з кешу без нового навчання. Бюджет пошуку вимірюється саме кількістю унікальних навчань.

Хромосоми одного покоління незалежні, тому навчаються паралельно в шести процесах, кожен з одним потоком обчислень. Мережа, що розійшлася під час навчання (прогнози містять NaN або $\infty$), отримує пристосованість $+\infty$ і відкидається відбором.

\begin{listing}[H]
\inputminted[firstline=187,lastline=198]{python}{../src/aiit/run_lab3.py}
\caption{Функція пристосованості (\texttt{src/aiit/run\_lab3.py})}
\label{lst:fitness}
\end{listing}

\section{Методика порівняння}

Сам по собі результат ГА мало що каже про ефективність алгоритму: покращення порівняно з ручним вибором могло б дати будь-яке систематичне дослідження простору. Тому ГА порівнюється з двома іншими методами за однакових умов.

\begin{itemize}
\item \textbf{Випадковий пошук} отримує рівно стільки навчань, скільки витратив ГА. Він використовує те саме зерно, тому перші 12 його хромосом збігаються з початковою популяцією ГА. Отже, методи відрізняються лише тим, що відбувається після спільного старту: еволюцією або подальшим випадковим вибором.
\item \textbf{«Наївний» ГА} --- той самий алгоритм із вимкненими відкиданням дублікатів та іммігрантами. Він показує, від чого захищають ці механізми.
\end{itemize}

Один запуск не дає змоги відрізнити кращий метод від удачливішого. Тому всі три методи запускалися з десяти різних зерен пошуку (42--51). Зерно навчання мереж при цьому лишалося незмінним, тож кожна хромосома мала ту саму пристосованість в усіх запусках, а змінювалася лише траєкторія пошуку.

\section{Результати запуску із зерном 42}

\subsection{Хід еволюції}

Програма запускається з кореня робочого простору командою \texttt{uv run --project <тека дисципліни> python <тека дисципліни>/lab3/run.py}. Нижче наведено фрагмент її консольного виводу: для кожного покоління --- найкраща і медіанна валідаційна RMSE, різноманітність, кількість іммігрантів і накопичена кількість унікальних навчань.

\begin{minted}[linenos=false]{text}
Search space: 18432 configurations
Generation 0: best=0.5573 median=0.7165 diversity=0.71 immigrants=0 evaluations=12
Generation 1: best=0.5548 median=0.5895 diversity=0.59 immigrants=0 evaluations=22
Generation 3: best=0.5468 median=0.5568 diversity=0.43 immigrants=0 evaluations=40
Generation 6: best=0.5394 median=0.5519 diversity=0.42 immigrants=0 evaluations=65
Generation 8: best=0.5368 median=0.5463 diversity=0.33 immigrants=0 evaluations=78
Generation 14: best=0.5368 median=0.5498 diversity=0.36 immigrants=0 evaluations=117
Seed 42, genetic: best RMSE(val)=0.5368 after 117 trainings, 58 s
Seed 42, naive-genetic: best RMSE(val)=0.5441 after 100 trainings, 48 s
Seed 42, random-search: best RMSE(val)=0.5438 after 117 trainings, 30 s
\end{minted}
\captionof{listing}{Фрагмент виводу програми (проміжні покоління пропущено)}
\label{lst:log}

\begin{figure}[H]
  \centering
  \includegraphics[width=\linewidth]{lab3/ga_convergence}
  \caption{Пристосованість і різноманітність популяції за поколіннями}
  \label{fig:convergence}
\end{figure}

Уже випадкова початкова популяція містила конфігурацію, кращу за модель лабораторної роботи №2 (0,557 проти 0,596). Це очікувано: модель ЛР2 обиралася вручну з чотирьох варіантів, а тут випадково оцінено 12 точок широкого простору. Далі найкраща RMSE монотонно знижувалася до 0,5368 у поколінні 8 і після цього не змінювалася. Медіана популяції за перші три покоління впала з 0,717 до 0,557: відбір швидко вилучив явно невдалі конфігурації, і популяція зосередилася в зоні хороших розв'язків.

Протягом 15 поколінь ГА навчив 117 унікальних мереж замість $12 \cdot 15 = 180$. Решта оцінок --- елітні особини і повторно знайдені хромосоми, пристосованість яких узято з кешу.

\subsection{Еволюція гіперпараметрів}

\begin{figure}[H]
  \centering
  \includegraphics[width=\linewidth]{lab3/gene_evolution}
  \caption{Частка популяції з кожним значенням кожного гена за поколіннями}
  \label{fig:genes}
\end{figure}

Рисунок~\ref{fig:genes} показує, до яких значень зійшлася популяція.

\begin{itemize}
\item \textbf{Ємність мережі зросла.} З покоління 2 більшість популяції становлять мережі з трьома шарами. Ширина 32 зникла вже в поколінні 2, а ширина 64 далі траплялася лише поодиноко.
\item \textbf{Регуляризацію відкинуто.} Після покоління 4 BatchNorm траплявся лише поодиноко, а в останніх поколіннях зник зовсім. Weight decay в останніх поколіннях набуває лише значень 0 і $10^{-5}$, Dropout --- переважно 0 або 0,1.
\item \textbf{Швидкість навчання зросла} до максимального значення $10^{-2}$, розмір батча стабілізувався на 128. Функцію активації GELU обрала більшість популяції, хоча на початку переважав LeakyReLU.
\end{itemize}

Ці результати незалежно підтверджують висновок лабораторної роботи №2. Там погіршення від регуляризації пояснювалося недонавчанням: даних у кілька разів більше, ніж параметрів мережі, тож обмежувати ємність немає потреби. ГА прийшов до того самого, не маючи жодної апріорної інформації: він збільшив ємність мережі в 16 разів (43\,521 параметр проти 2\,689) і відмовився від регуляризації.

Два гени зійшлися до крайніх значень діапазону: ширина 256 і швидкість навчання $10^{-2}$. Це ознака того, що оптимум може лежати за межами заданого простору пошуку, тож наступним кроком варто було б розширити ці діапазони. Висока швидкість навчання має й іншу причину: пристосованість вимірюється після фіксованих 20 епох, і ГА оптимізує саме цю величину. Тобто він знаходить конфігурацію, яка найшвидше навчається за 20 епох, а не ту, що досягла б найкращого результату за необмеженого навчання.

\subsection{Порівняння з випадковим пошуком}

\begin{figure}[H]
  \centering
  \includegraphics[width=0.9\linewidth]{lab3/search_comparison}
  \caption{Найкраща знайдена RMSE залежно від кількості навчань (точки --- окремі конфігурації)}
  \label{fig:search}
\end{figure}

Усі три методи стартували з однієї й тієї ж популяції, тому перші 12 точок кривих збігаються. Приблизно з 50-го навчання ГА лишався попереду обох інших методів і завершив пошук із RMSE 0,5368 проти 0,5438 у випадкового пошуку та 0,5441 у наївного ГА. Видно і якісну відмінність: точки випадкового пошуку рівномірно розкидані до кінця, тоді як точки ГА з часом зосереджуються біля найкращого значення. ГА витрачає бюджет на дослідження околу хороших розв'язків, а не на випадкові конфігурації.

\begin{table}[H]
\caption{Фінальні конфігурації (зерно 42)}
\label{tab:finals}
\begin{tabularx}{\linewidth}{|X|c|r|r|r|r|}
\hline
\textbf{Метод} & \textbf{Шари} & \textbf{Парам.} & \textbf{RMSE val} & \textbf{RMSE test} & \textbf{$R^2$ test} \\
\hline
ЛР2: baseline-relu & 64--32 & 2\,689 & 0,5964 & 0,6087 & 0,7186 \\
\hline
Випадковий пошук & 256--128--64 & 43\,521 & 0,5438 & \textbf{0,5380} & \textbf{0,7802} \\
\hline
Наївний ГА & 256--128 & 35\,329 & 0,5441 & 0,5463 & 0,7734 \\
\hline
Генетичний алгоритм & 256--128--64 & 43\,521 & \textbf{0,5368} & 0,5511 & 0,7694 \\
\hline
\end{tabularx}
\end{table}

\begin{table}[H]
\caption{Гіперпараметри фінальних конфігурацій}
\label{tab:finalparams}
\begin{tabularx}{\linewidth}{|X|c|c|c|c|c|c|}
\hline
\textbf{Метод} & \textbf{Актив.} & \textbf{Dropout} & \textbf{BN} & \textbf{L2} & \textbf{LR} & \textbf{Батч} \\
\hline
ЛР2: baseline-relu & ReLU & 0 & ні & 0 & $10^{-3}$ & 256 \\
\hline
Випадковий пошук & GELU & 0 & ні & $10^{-5}$ & $3 \cdot 10^{-3}$ & 64 \\
\hline
Наївний ГА & ReLU & 0 & ні & $10^{-4}$ & $10^{-2}$ & 128 \\
\hline
Генетичний алгоритм & GELU & 0,1 & ні & 0 & $10^{-2}$ & 128 \\
\hline
\end{tabularx}
\end{table}

Усі три методи покращили модель ЛР2 на 9--12~\% за тестовою RMSE: з 0,609 до 0,538--0,551. На валідаційній вибірці найкращим є ГА, але на тестовій --- випадковий пошук. Різниця між методами на тесті (до 0,013) того самого порядку, що й на валідації (0,007), тому за одним запуском визначити кращий метод неможливо. Саме для цього виконано серію запусків із різними зернами (розділ~\ref{sec:seeds}).

Тестова RMSE моделі ЛР2 тут становить 0,6087, а не 0,6165, як у лабораторній роботі №2. Причина в тому, що для паралельного навчання цикл навчання переписано з однопотоковими обчисленнями і власним перемішуванням батчів, тож ваги навчаються в іншому порядку. Для коректного порівняння модель ЛР2 перенавчено тією самою процедурою, що й усі кандидати пошуку.

\begin{figure}[H]
  \centering
  \includegraphics[width=0.62\linewidth]{lab3/pred_vs_actual}
  \caption{Прогнозовані значення проти фактичних на тестовій вибірці (найкраща конфігурація ГА)}
  \label{fig:pred}
\end{figure}

\subsection{Виродження популяції}

Права частина рис.~\ref{fig:convergence} порівнює різноманітність популяції в ГА і в наївному ГА. Покоління 0--2 в обох алгоритмах однакові: вони отримали ті самі випадкові числа, а дублікатів ще не виникало. З покоління 3, коли ГА вперше відкинув дублікат, траєкторії розходяться, а після покоління 6 різниця стає суттєвою. У наївному ГА різноманітність падає до 0,15 у поколінні 9: на той момент у популяції лишилося лише 4 різні хромосоми з 12, а медіанна RMSE збіглася з найкращою (0,5466), тобто більшу частину популяції становили копії лідера. Це і є виродження: схрещування копій дає ті самі копії, і з покоління 8 по 9 наївний ГА навчив лише дві нові мережі.

У ГА із захистом різноманітність не опускалася нижче 0,32, а всі 12 особин кожного покоління були унікальними. Порогу 0,3 вона жодного разу не досягла, тож іммігранти в цьому запуску не знадобилися: від виродження вберегло саме відкидання дублікатів. Наївний ГА, що застряг у поколіннях 5--9, витратив на 17 навчань менше, але знайшов гірший розв'язок (0,5441 проти 0,5368).

\section{Повторні запуски з різними зернами}
\label{sec:seeds}

Кожен із трьох методів запущено з десяти зерен пошуку (42--51). Для кожного запуску переможця пошуку перенавчено і виміряно на тестовій вибірці. Підсумки наведено в табл.~\ref{tab:seeds} і на рис.~\ref{fig:seeds}.

\begin{table}[H]
\caption{Підсумки 10 запусків кожного методу (середнє $\pm$ стандартне відхилення)}
\label{tab:seeds}
\begin{tabularx}{\linewidth}{|X|c|c|c|}
\hline
\textbf{Метод} & \textbf{Навчань} & \textbf{RMSE val} & \textbf{RMSE test} \\
\hline
Генетичний алгоритм & 119,7 & $\mathbf{0{,}5436 \pm 0{,}0041}$ & $0{,}5502 \pm 0{,}0071$ \\
\hline
Наївний ГА & 94,3 & $0{,}5464 \pm 0{,}0046$ & $\mathbf{0{,}5499 \pm 0{,}0079}$ \\
\hline
Випадковий пошук & 119,7 & $0{,}5516 \pm 0{,}0051$ & $0{,}5545 \pm 0{,}0080$ \\
\hline
ЛР2: baseline-relu & --- & 0,5964 & 0,6087 \\
\hline
\end{tabularx}
\end{table}

\begin{figure}[H]
  \centering
  \includegraphics[width=\linewidth]{lab3/seed_study}
  \caption{Порівняння методів за 10 запусками: середня крива пошуку (ліворуч) і RMSE переможців кожного запуску (праворуч)}
  \label{fig:seeds}
\end{figure}

\subsection{Генетичний алгоритм проти випадкового пошуку}

\textbf{На валідаційній вибірці}, тобто за функцією пристосованості, ГА виявився кращим за випадковий пошук \textbf{в усіх 10 запусках} із тим самим бюджетом навчань. Середня перевага становить 0,0080, і за парним t-критерієм вона статистично значуща ($t = 6{,}70$, $p = 0{,}0001$). Отже, як метод оптимізації заданої функції ГА однозначно ефективніший. Найкращу відому конфігурацію (RMSE 0,5368) ГА знайшов у двох запусках різними траєкторіями (зерна 42 і 46), а випадковий пошук жодного разу не опустився нижче 0,5438.

\textbf{На тестовій вибірці} перевага скорочується: ГА кращий у 7 запусках із 10, середня різниця 0,0044, і вона статистично не значуща ($t = 1{,}17$, $p = 0{,}27$). Причину видно з різниці між тестовою і валідаційною RMSE переможців: у ГА вона в середньому становить $+0{,}0065$, у випадкового пошуку --- $+0{,}0029$.

Це прояв так званого «прокляття переможця». Валідаційна RMSE сама є оцінкою, обчисленою на скінченній вибірці з 4\,128 записів, і містить випадкову похибку. ГА зосереджує навчання в околі найкращого розв'язку, тож порівнює між собою багато майже однакових конфігурацій. Серед них перемагає та, для якої похибка валідаційної оцінки виявилася найсприятливішою. ГА оптимізує саме ту функцію, яку йому задали, разом із її шумом. Практичний висновок: якщо ГА використовують для підбору гіперпараметрів, функцію пристосованості варто робити стійкішою до шуму, наприклад обчислювати RMSE крос-валідацією або на більшій валідаційній вибірці, а фінальний результат обов'язково перевіряти на окремій тестовій вибірці, як це зроблено тут.

Окремо слід зазначити витрати часу. За однакової кількості навчань ГА працював приблизно вдвічі довше: 58 с проти 30 с для зерна 42. Випадковий пошук відразу передає всі конфігурації шести паралельним процесам, тоді як ГА мусить дочекатися оцінки всього покоління, перш ніж формувати наступне, тож частина процесів простоює.

\subsection{Ефект захисту від виродження}

У всіх 10 запусках наївного ГА популяція вироджувалася: мінімальна різноманітність у середньому становила 0,19 і опускалася до 0,15. ГА із захистом не опускався нижче 0,31. Іммігранти знадобилися лише двічі за 10 запусків (зерна 47 і 51, по одному разу), тобто основну роботу виконало відкидання дублікатів.

Водночас ліва частина рис.~\ref{fig:seeds} показує, що за однакової кількості навчань наївний ГА не гірший. На перших 40 навчаннях він навіть попереду (середня найкраща RMSE 0,5525 проти 0,5565), бо без відкидання дублікатів швидше концентрується біля лідера. ГА із захистом наздоганяє його лише після 70 навчань (0,5471 проти 0,5475). Різниця між ними на валідації погранично значуща ($p = 0{,}07$ за t-критерієм), а на тесті відсутня ($p = 0{,}89$).

Основний наслідок виродження тут --- не гірший розв'язок, а марнування поколінь. Виродженна популяція складається з копій уже оцінених хромосом, і за ті самі 15 поколінь наївний ГА навчив у середньому лише 94 нові мережі проти 120. Пошук фактично зупиняється раніше, ніж закінчується відведена кількість поколінь.

\subsection{Стабільність знайдених гіперпараметрів}

Переможці різних запусків ГА відрізняються, але мають спільні риси. Жоден із 30 переможців (10 запусків $\times$ 3 методи) не використовує BatchNorm. У 9 з 10 запусків ГА переможець має ширину першого шару 256, у 8 --- Dropout 0, у всіх --- швидкість навчання $3 \cdot 10^{-3}$ або $10^{-2}$. Тобто картина, яку показує рис.~\ref{fig:genes}, не випадкова: висока ємність і відмова від регуляризації стабільно виграють на цій задачі.

\section{Висновки}

У ході лабораторної роботи реалізовано генетичний алгоритм для підбору гіперпараметрів нейронної мережі та досліджено його ефективність порівняно з випадковим пошуком.

\begin{enumerate}
\item Реалізовано власний генетичний алгоритм із турнірним відбором, рівномірним схрещуванням, мутацією з урахуванням типу гена та елітизмом. Від виродження популяції його захищають відкидання дублікатів і введення випадкових іммігрантів. Алгоритм досліджує простір із 18\,432 конфігурацій восьми гіперпараметрів, витрачаючи близько 120 навчань, тобто менш ніж 1~\% простору.

\item ГА знайшов конфігурацію (три шари $256 \to 128 \to 64$, GELU, Dropout 0,1, швидкість навчання $10^{-2}$) з тестовою RMSE 0,551 і $R^2 = 0{,}769$, що на 9,5~\% краще за модель лабораторної роботи №2 (RMSE 0,609). За 10 запусками середня тестова RMSE переможця ГА становить $0{,}550 \pm 0{,}007$.

\item За функцією пристосованості ГА перевершив випадковий пошук з тим самим бюджетом в усіх 10 запусках, і ця перевага статистично значуща ($p = 0{,}0001$). На тестовій вибірці перевага зменшилася до 0,004 і стала незначущою ($p = 0{,}27$). Причина --- ефект «прокляття переможця»: ГА оптимізує валідаційну оцінку разом із її випадковою похибкою. Для практичного застосування функцію пристосованості слід робити стійкішою до шуму, наприклад через крос-валідацію.

\item Без захисту від виродження різноманітність популяції падала до 0,15, і популяцію заповнювали копії лідера. Захищений ГА утримував різноманітність не нижче 0,31. За однакової кількості навчань наївний ГА не поступався якістю, але виродженна популяція марнувала покоління: за ті самі 15 поколінь він навчив приблизно на п'яту частину менше нових мереж.

\item ГА незалежно підтвердив висновок лабораторної роботи №2 про недонавчання моделі: в усіх запусках він відмовлявся від BatchNorm, мінімізував регуляризацію і збільшував ємність мережі (до 16 разів за кількістю параметрів). Ширина шару і швидкість навчання зійшлися до меж заданих діапазонів, тому простір пошуку варто розширити.

\item Результат ГА визначається функцією пристосованості. Оскільки пристосованість вимірюється після фіксованих 20 епох, ГА відбирає конфігурації, що найшвидше навчаються, зокрема з максимальною швидкістю навчання. Вибір функції пристосованості є ключовим рішенням під час застосування еволюційних алгоритмів.
\end{enumerate}

\end{document}
